\documentclass[11pt]{article}

\usepackage[a4paper, margin=1in]{geometry}
\usepackage{graphicx}
\usepackage{amsmath}
\usepackage{booktabs}
\usepackage{caption}
\usepackage{float}
\usepackage{hyperref}
\usepackage{array}
\usepackage{enumitem}
\usepackage[most]{tcolorbox}

\hypersetup{colorlinks=true, urlcolor=blue, linkcolor=black}

\title{\vspace{-2cm}}
\date{}

\begin{document}

% ============================================================
% HEADER / COVER BLOCK (matches the Experiment 1 report style)
% ============================================================
\begin{center}
    {\Large \textbf{Shiv Nadar University Chennai}}\\[4pt]
\end{center}

\noindent
\begin{tabular}{@{}p{4.2cm}l@{}}
    Degree \& Branch    & B.Tech Artificial Intelligence \& Data Science, Semester V \\
    Subject             & CS3807 -- Deep Learning Laboratory \\
    Name                & Madhav.S \\
    Section             & AIDS `B' \\
    Register Number     & 24011101066 \\
\end{tabular}

\vspace{0.6cm}
\begin{center}
    {\Large \textbf{Experiment 2}}\\[6pt]
    {\large \textbf{Implementation of a Multi-Layer Perceptron (MLP) for Multi-Class Image Classification}}
\end{center}

\vspace{0.4cm}
\hrule
\vspace{0.6cm}

% ============================================================
% 1. OBJECTIVE
% ============================================================
\section{Objective}
The goal of this experiment was to implement a Multi-Layer Perceptron (MLP) using
TensorFlow/Keras on the Fashion-MNIST dataset, covering image preprocessing,
flattening, model construction, training, evaluation, and automated hyperparameter
optimization.

% ============================================================
% 2. LEARNING OUTCOMES
% ============================================================
\section{Learning Outcomes}
Through this experiment the following outcomes were achieved:
\begin{enumerate}[label=\arabic*.]
    \item Explained the architecture of an MLP.
    \item Preprocessed an image dataset for a dense-layer network.
    \item Built and trained an MLP using Keras.
    \item Evaluated classification performance using multiple metrics.
    \item Performed automated hyperparameter optimization.
    \item Recommended the best model configuration found by the search.
\end{enumerate}

% ============================================================
% 3. BACKGROUND THEORY
% ============================================================
\section{Background Theory}
An MLP consists of an input layer, one or more hidden layers, and an output layer.
Each layer $l$ computes a weighted sum of its inputs, adds a bias, and passes the
result through an activation function:
\begin{equation}
    z^{(l)} = W^{(l)} a^{(l-1)} + b^{(l)}, \qquad a^{(l)} = f\left(z^{(l)}\right)
\end{equation}

The output layer uses a softmax activation to produce a probability distribution
over the 10 classes:
\begin{equation}
    \hat{y} = \mathrm{Softmax}(z)
\end{equation}

Training minimizes the categorical cross-entropy loss:
\begin{equation}
    \mathcal{L} = -\sum_i y_i \log(\hat{y}_i)
\end{equation}

The recommended architecture used for the baseline model in this experiment was:
\begin{center}
    $784 \rightarrow \text{Dense}(128, \text{ReLU}) \rightarrow \text{Dense}(64, \text{ReLU}) \rightarrow \text{Dense}(10, \text{Softmax})$
\end{center}

% ============================================================
% 4. DATASET
% ============================================================
\section{Dataset}
The Fashion-MNIST dataset was used for this experiment.

\begin{table}[H]
\centering
\caption{Dataset Summary (verified at load time)}
\begin{tabular}{@{}ll@{}}
\toprule
\textbf{Property} & \textbf{Value} \\
\midrule
Training Images & 60{,}000, shape $(60000, 28, 28)$ \\
Testing Images  & 10{,}000, shape $(10000, 28, 28)$ \\
Classes         & 10 \\
Image Size      & $28 \times 28$ \\
Pixel Range     & 0 to 255 \\
\bottomrule
\end{tabular}
\end{table}

The 10 classes are: T-shirt/top, Trouser, Pullover, Dress, Coat, Sandal, Shirt,
Sneaker, Bag, and Ankle boot.

% ============================================================
% 5. IMAGE PREPROCESSING
% ============================================================
\section{Image Preprocessing}
Fashion-MNIST images are stored as $28 \times 28$ matrices. Dense layers require a
one-dimensional feature vector, so each image was flattened:
\begin{equation}
    28 \times 28 = 784
\end{equation}

The following preprocessing steps were performed:
\begin{enumerate}[label=\arabic*.]
    \item Loaded Fashion-MNIST.
    \item Displayed sample images (Figure~\ref{fig:sample_images}).
    \item Flattened images from $(N, 28, 28)$ to $(N, 784)$.
    \item Normalized pixel values from $[0, 255]$ to $[0, 1]$.
    \item Converted labels into one-hot vectors.
\end{enumerate}

\begin{table}[H]
\centering
\caption{Tensor Shapes Before and After Preprocessing}
\begin{tabular}{@{}lll@{}}
\toprule
\textbf{Tensor} & \textbf{Before} & \textbf{After} \\
\midrule
Training Images & $(60000, 28, 28)$ & $(60000, 784)$ \\
Testing Images  & $(10000, 28, 28)$ & $(10000, 784)$ \\
Training Labels & $(60000,)$        & $(60000, 10)$ one-hot \\
Testing Labels  & $(10000,)$        & $(10000, 10)$ one-hot \\
\bottomrule
\end{tabular}
\end{table}

% ============================================================
% 6. EXPERIMENTAL PROCEDURE
% ============================================================
\section{Experimental Procedure}

\subsection*{Task 1: Dataset Exploration}
The dataset dimensions were printed and ten sample images were displayed
(Figure~\ref{fig:sample_images}), along with the class distribution across the
training set (Figure~\ref{fig:class_distribution}).

\subsection*{Task 2: Data Preprocessing}
Images were flattened and normalized as described in Section~5, with tensor
shapes printed before and after each step (Table~2).

\subsection*{Task 3: Model Construction}
The baseline MLP was constructed with the recommended architecture:\\
Input(784) $\rightarrow$ Dense(128, ReLU) $\rightarrow$ Dense(64, ReLU) $\rightarrow$ Dense(10, Softmax).

\begin{table}[H]
\centering
\caption{Baseline Model Summary}
\begin{tabular}{@{}lll@{}}
\toprule
\textbf{Layer} & \textbf{Output Shape} & \textbf{Parameters} \\
\midrule
Dense (ReLU)     & $(\text{None}, 128)$ & 100{,}480 \\
Dense\_1 (ReLU)  & $(\text{None}, 64)$  & 8{,}256 \\
Dense\_2 (Softmax) & $(\text{None}, 10)$ & 650 \\
\midrule
\multicolumn{2}{@{}l}{Total Parameters}    & 109{,}386 \\
\multicolumn{2}{@{}l}{Trainable Parameters} & 109{,}386 \\
\multicolumn{2}{@{}l}{Non-trainable Parameters} & 0 \\
\bottomrule
\end{tabular}
\end{table}

\subsection*{Task 4: Model Training}
The baseline model was compiled with the Adam optimizer (learning rate 0.001),
categorical cross-entropy loss, and accuracy as the evaluation metric. It was
trained for 20 epochs with a batch size of 32, using a 10\% validation split.

\begin{table}[H]
\centering
\caption{Baseline Training Log (Selected Epochs)}
\begin{tabular}{@{}ccccc@{}}
\toprule
\textbf{Epoch} & \textbf{Train Acc.} & \textbf{Train Loss} & \textbf{Val. Acc.} & \textbf{Val. Loss} \\
\midrule
1  & 0.8189 & 0.5036 & 0.8470 & 0.4107 \\
5  & 0.8907 & 0.2924 & 0.8813 & 0.3264 \\
10 & 0.9105 & 0.2377 & 0.8835 & 0.3246 \\
15 & 0.9230 & 0.2032 & 0.8907 & 0.3336 \\
20 & 0.9337 & 0.1748 & 0.8847 & 0.3890 \\
\bottomrule
\end{tabular}
\end{table}

\subsection*{Task 5: Model Evaluation}
The trained baseline model was evaluated on the held-out test set.

\begin{table}[H]
\centering
\caption{Baseline Model Test Performance}
\begin{tabular}{@{}ll@{}}
\toprule
\textbf{Metric} & \textbf{Value} \\
\midrule
Accuracy  & 0.8836 \\
Precision (macro) & 0.8853 \\
Recall (macro)    & 0.8836 \\
F1-score (macro)  & 0.8839 \\
\bottomrule
\end{tabular}
\end{table}

\begin{table}[H]
\centering
\caption{Baseline Model -- Per-Class Classification Report}
\begin{tabular}{@{}lccc@{}}
\toprule
\textbf{Class} & \textbf{Precision} & \textbf{Recall} & \textbf{F1-score} \\
\midrule
T-shirt/top & 0.84 & 0.81 & 0.83 \\
Trouser     & 0.99 & 0.97 & 0.98 \\
Pullover    & 0.76 & 0.84 & 0.80 \\
Dress       & 0.91 & 0.89 & 0.90 \\
Coat        & 0.84 & 0.75 & 0.79 \\
Sandal      & 0.94 & 0.98 & 0.96 \\
Shirt       & 0.69 & 0.72 & 0.70 \\
Sneaker     & 0.96 & 0.94 & 0.95 \\
Bag         & 0.96 & 0.98 & 0.97 \\
Ankle boot  & 0.96 & 0.94 & 0.95 \\
\midrule
\textbf{Macro avg}    & 0.89 & 0.88 & 0.88 \\
\textbf{Weighted avg} & 0.89 & 0.88 & 0.88 \\
\bottomrule
\end{tabular}
\end{table}

% ============================================================
% 7. HYPERPARAMETER OPTIMIZATION
% ============================================================
\section{Hyperparameter Optimization}
Instead of manually selecting hyperparameters, automated hyperparameter
optimization was performed using \texttt{RandomizedSearchCV} together with the
SciKeras wrapper (\texttt{KerasClassifier}), with 5-fold cross-validation.

\begin{table}[H]
\centering
\caption{Hyperparameter Search Space}
\begin{tabular}{@{}ll@{}}
\toprule
\textbf{Hyperparameter} & \textbf{Candidate Values} \\
\midrule
Hidden Layers        & 1, 2, 3 \\
Hidden Neurons       & 32, 64, 128, 256 \\
Learning Rate        & 0.1, 0.01, 0.001 \\
Batch Size           & 16, 32, 64, 128 \\
Epochs               & 10, 20, 30 \\
Optimizer            & SGD, Adam, RMSProp \\
Activation Function  & ReLU, Tanh, Sigmoid \\
Dropout Rate         & 0.0, 0.2, 0.5 \\
\bottomrule
\end{tabular}
\end{table}

The search was run on a stratified 10{,}000-sample subsample of the training set
(15 candidates $\times$ 5 folds = 75 model fits) to keep the search
computationally feasible, with the final optimized model retrained on the full
60{,}000-image training set using the best hyperparameters found.

% ============================================================
% 8. MANDATORY PLOTS
% ============================================================
\section{Mandatory Plots}

\begin{figure}[H]
    \centering
    \includegraphics[width=1\textwidth]{sample_images.pdf}
    \caption{Sample Fashion-MNIST Images}
    \label{fig:sample_images}
\end{figure}
\noindent\textbf{Inference:} Each of the ten classes is visually distinct even
at $28 \times 28$ resolution, though some categories (e.g.\ Shirt, T-shirt/top,
Pullover, and Coat) share similar silhouettes.

\begin{figure}[H]
    \centering
    \includegraphics[width=1.1\textwidth]{ class_distribution.pdf}
    \caption{Class Distribution (Training Set)}
    \label{fig:class_distribution}
\end{figure}
\noindent\textbf{Inference:} All 10 classes are represented by exactly 6{,}000
training images each, so the dataset is perfectly balanced. This means accuracy
is a reliable metric here and no class-imbalance handling was required.

\begin{figure}[H]
    \centering
    \includegraphics[width=0.85\textwidth]{ training_accuracy_vs_epoch.pdf}
    \caption{Training Accuracy vs Epoch}
    \label{fig:train_acc}
\end{figure}
\noindent\textbf{Inference:} Training accuracy increases steadily and
monotonically from 81.9\% to 93.4\% over 20 epochs, showing the model
continues to fit that the training data gets better at every epoch.

\begin{figure}[H]
    \centering
    \includegraphics[width=0.85\textwidth]{ validation_accuracy_vs_epoch.pdf}
    \caption{Validation Accuracy vs Epoch}
    \label{fig:val_acc}
\end{figure}
\noindent\textbf{Inference:} Validation accuracy rises quickly in the first
few epochs and then plateaus around 88--89\%, well below the training accuracy
of 93\%+ by epoch 20. This shows that the model begins to overfit on the training set in the later epochs.

\begin{figure}[H]
    \centering
    \includegraphics[width=0.85\textwidth]{ training_loss_vs_epoch.pdf}
    \caption{Training Loss vs Epoch}
    \label{fig:train_loss}
\end{figure}
\noindent\textbf{Inference:} Training loss decreases smoothly and consistently
from 0.50 to 0.17 across all 20 epochs, almost similar to the training accuracy curve. This confirms stable optimization with the Adam optimizer.

\begin{figure}[H]
    \centering
    \includegraphics[width=0.85\textwidth]{ validation_loss_vs_epoch.pdf}
    \caption{Validation Loss vs Epoch}
    \label{fig:val_loss}
\end{figure}
\noindent\textbf{Inference:} Validation loss decreases until around epoch
10--15 and then begins to increase(ending near 0.39), while training
loss keeps decreasing. This is the classic divergence pattern that confirms overfitting
in the later epochs.

\begin{figure}[H]
    \centering
    \includegraphics[width=0.75\textwidth]{ confusion_matrix.pdf}
    \caption{Confusion Matrix -- Baseline Model}
    \label{fig:confusion_baseline}
\end{figure}
\noindent\textbf{Inference:} Most confusion occurs between Shirt, T-shirt/top,
Pullover, and Coat -- all upper-body garments with similar silhouettes at low
resolution. Trouser, Sandal, Bag, and Ankle boot are classified almost
perfectly, consistent with their more unique shapes.

\begin{figure}[H]
    \centering
    \includegraphics[width=0.85\textwidth]{ hyperparameter_search_results.pdf}
    \caption{Hyperparameter Search Results}
    \label{fig:search_results}
\end{figure}
\noindent\textbf{Inference:} Cross-validation accuracy varies noticeably across
the 15 sampled candidates, ranging from weak configurations up to the best candidate at 85.79\% mean CV accuracy, showing that the hyperparameter choice has a real, measurable effect on the model quality even on a subsampled
training set.

\begin{figure}[H]
    \centering
    \includegraphics[width=0.75\textwidth]{ best_model_comparison.pdf}
    \caption{Best Model Accuracy Comparison}
    \label{fig:best_model_comparison}
\end{figure}
\noindent\textbf{Inference:} The optimized model outperforms the baseline on
every metric (accuracy, precision, recall, and F1-score), though the margin is
slight (under 0.4 percentage points), indicating that the baseline architecture was
already a reasonably strong choice for this dataset.

\begin{figure}[H]
    \centering
    \includegraphics[width=0.75\textwidth]{ confusion_matrix_optimized.pdf}
    \caption{Confusion Matrix -- Optimized Model (bonus)}
    \label{fig:confusion_optimized}
\end{figure}
\noindent\textbf{Inference:} The optimized model shows the same overall
confusion pattern as the baseline (Shirt/T-shirt/Pullover/Coat remain the
hardest to separate), but with slightly fewer total misclassifications,
consistent with its marginally higher test accuracy.

% ============================================================
% SOURCE CODE
% ============================================================
\section{Source Code}
The complete implementation, exactly as run in Colab, is hosted on GitHub.
\\
\url{https://github.com/smadhav2024/Deep_Learning_Lab/tree/main/Lab%202}
\\
The repository contains the full notebook \texttt{deep\_learning\_lab\_2.ipynb},
the \texttt{requirements.txt} needed to run it, the raw \texttt{.eps} plots,
and this report's Latex code.

% ============================================================
% 9. RESULTS
% ============================================================
\section{Results}

\begin{table}[H]
\centering
\caption{Best Hyperparameters Found (RandomizedSearchCV, 5-fold CV)}
\begin{tabular}{@{}ll@{}}
\toprule
\textbf{Hyperparameter} & \textbf{Best Value} \\
\midrule
Hidden Layers       & 1 \\
Hidden Neurons      & 128 \\
Learning Rate       & 0.001 \\
Batch Size          & 128 \\
Optimizer           & Adam \\
Activation Function & Sigmoid \\
Epochs              & 30 \\
Dropout             & 0.0 \\
\midrule
Cross-validation Accuracy & 0.8579 \\
Testing Accuracy (after full retrain) & 0.8871 \\
\bottomrule
\end{tabular}
\end{table}

\begin{table}[H]
\centering
\caption{Performance Comparison: Baseline vs Optimized}
\begin{tabular}{@{}lcccc@{}}
\toprule
\textbf{Metric} & \textbf{Baseline} & \textbf{Optimized} \\
\midrule
Accuracy  & 0.8836 & 0.8871 \\
Precision & 0.8853 & 0.8871 \\
Recall    & 0.8836 & 0.8871 \\
F1-score  & 0.8839 & 0.8868 \\
\bottomrule
\end{tabular}
\end{table}

% ============================================================
% ADDITIONAL TASK (K4): PERCEPTRON LEARNING ALGORITHM FOR XOR
% (placed after Results, before Discussion, as requested)
% ============================================================
\section{Additional Task (K4): Perceptron Learning Algorithm for the XOR Gate}

\subsection*{Objective}
The perceptron learning algorithm was implemented from scratch in Python for
the XOR logic gate. The weights were logged after every update, the decision
boundary was plotted after each weight update, and the pattern that prevents
the single-layer perceptron from converging on XOR was analyzed.

\subsection*{Implementation}
A single-layer perceptron with two input weights and a bias was trained on the
XOR truth table using the standard perceptron update rule:
\begin{equation}
    w \leftarrow w + \eta\,(t - \hat{y})\,x, \qquad b \leftarrow b + \eta\,(t - \hat{y})
\end{equation}
where $\hat{y} = \mathrm{step}(w \cdot x + b)$, $t$ is the target output, and
$\eta = 0.1$ is the learning rate. Weights and bias were initialized to zero
and the perceptron was trained for 5 epochs over the 4 XOR samples
$\{(0,0){\to}0,\ (0,1){\to}1,\ (1,0){\to}1,\ (1,1){\to}0\}$, with a weight
update logged and the weight/bias state recorded every time a sample was
misclassified.

\subsection*{Weights After Each Update}

\begin{table}[H]
\centering
\small
\caption{Perceptron Weight Updates on the XOR Gate ($\eta = 0.1$, initial $w = [0,0]$, $b = 0$)}
\begin{tabular}{@{}cccccccc@{}}
\toprule
\textbf{Update} & \textbf{Epoch} & \textbf{Input} & \textbf{Target} & \textbf{Predicted} & \textbf{$w_1$} & \textbf{$w_2$} & \textbf{Bias} \\
\midrule
1  & 1 & (0,0) & 0 & 1 & 0.0000  & 0.0000 & $-0.1000$ \\
2  & 1 & (0,1) & 1 & 0 & 0.0000  & 0.1000 & 0.0000 \\
3  & 1 & (1,1) & 0 & 1 & $-0.1000$ & 0.1000 & $-0.1000$ \\
4  & 2 & (0,1) & 1 & 0 & $-0.1000$ & 0.1000 & 0.0000 \\
5  & 2 & (1,0) & 1 & 0 & 0.0000  & 0.1000 & 0.1000 \\
6  & 2 & (1,1) & 0 & 1 & $-0.1000$ & 0.1000 & 0.0000 \\
7  & 3 & (0,0) & 0 & 1 & $-0.1000$ & 0.1000 & $-0.1000$ \\
8  & 3 & (0,1) & 1 & 0 & $-0.1000$ & 0.1000 & 0.0000 \\
9  & 3 & (1,0) & 1 & 0 & 0.0000  & 0.1000 & 0.1000 \\
10 & 3 & (1,1) & 0 & 1 & $-0.1000$ & 0.1000 & 0.0000 \\
11 & 4 & (0,0) & 0 & 1 & $-0.1000$ & 0.1000 & $-0.1000$ \\
12 & 4 & (0,1) & 1 & 0 & $-0.1000$ & 0.1000 & 0.0000 \\
13 & 4 & (1,0) & 1 & 0 & 0.0000  & 0.1000 & 0.1000 \\
14 & 4 & (1,1) & 0 & 1 & $-0.1000$ & 0.1000 & 0.0000 \\
15 & 5 & (0,0) & 0 & 1 & $-0.1000$ & 0.1000 & $-0.1000$ \\
16 & 5 & (0,1) & 1 & 0 & $-0.1000$ & 0.1000 & 0.0000 \\
17 & 5 & (1,0) & 1 & 0 & 0.0000  & 0.1000 & 0.1000 \\
18 & 5 & (1,1) & 0 & 1 & $-0.1000$ & 0.1000 & 0.0000 \\
\bottomrule
\end{tabular}
\end{table}

\noindent Note that from Update~4 onward the weight/bias state simply repeats
the same 4-update cycle
$(-0.1,0.1,0.0) \to (0.0,0.1,0.1) \to (-0.1,0.1,0.0) \to (-0.1,0.1,-0.1) \to \dots$
every epoch, and the number of misclassifications per epoch was
$[3,\,3,\,4,\,4,\,4]$ over the 5 epochs -- it never drops to 0, i.e.\ the
perceptron never converges.

\subsection*{Decision Boundary Evolution}

\begin{figure}[H]
    \centering
    \includegraphics[width=1.1\textwidth]{xor_decision_boundary_evolution.eps}
    \caption{XOR Decision Boundary After Each of the 18 Weight Updates (Update 0 = initial state)}
    \label{fig:xor_boundary}
\end{figure}
\noindent\textbf{Inference:} The decision boundary line visibly rotates and
shifts after every update, moving to correct whichever point was just
misclassified. However, no single straight-line boundary manages to
place all the four XOR points on their correct sides simultaneously. The
boundary keeps going after the errors from one update to the next without ever
settling.

\begin{figure}[H]
    \centering
    \includegraphics[width=0.85\textwidth]{xor_training_error.eps}
    \caption{Perceptron Trained on XOR -- Non-Convergence}
    \label{fig:xor_error}
\end{figure}
\noindent\textbf{Inference:} The number of misclassified samples per epoch
($3, 3, 4, 4, 4$) never reaches zero and does not trend downward -- after
epoch 2 it stabilizes at 4 out of 4 samples misclassified at every epoch. This confirms that the perceptron has entered a permanent error cycle rather than
converging.

\begin{figure}[H]
    \centering
    \includegraphics[width=0.85\textwidth]{xor_weight_evolution.eps}
    \caption{XOR Weight and Bias Evolution Over Updates}
    \label{fig:xor_weights}
\end{figure}
\noindent\textbf{Inference:} After the initial updates (Updates 1--3), both
weights and the bias settle into a repeating oscillation between a small set
of fixed values instead of converging to a stable point. The direct numerical
evidence of the cyclical, non-convergent behaviour can be seen in
Figure~\ref{fig:xor_error}.

\subsection*{Analysis: Why the Perceptron Never Converges on XOR}
The pattern responsible for this non-convergence is that \textbf{XOR is not
linearly separable}. The two positive samples $(0,1)$ and $(1,0)$ and the two
negative samples $(0,0)$ and $(1,1)$ cannot be separated by any single
straight line in the 2D input space -- they form two pairs of diagonally
opposite corners of the unit square, so any line that correctly separates one
pair of opposite corners necessarily misclassifies the other pair. Every
weight update in Figure~\ref{fig:xor_boundary} rotates or shifts the boundary
to fix whichever point was just misclassified, but doing so always breaks a
different point that was previously correct, since no orientation of a single
straight line can keep all four points on their correct sides at once. This
is exactly why solving XOR requires at least one hidden layer, combining multiple linear boundaries
produces a non-linear decision region that a single-layer perceptron, limited
to one straight-line boundary, can never represent.

% ============================================================
% ADDITIONAL TASK (K5): MULTI-LAYER PERCEPTRON FOR XOR
% (placed immediately after the single-layer perceptron task, as requested)
% ============================================================
\section{Additional Task (K5): Multi-Layer Perceptron for the XOR Gate}

\begin{tcolorbox}[
    colback=white,
    colframe=black,
    boxrule=0.8pt,
    arc=1pt,
    left=6pt, right=6pt, top=4pt, bottom=4pt,
    title={\textbf{Task Summary}},
    fonttitle=\bfseries,
    colbacktitle=white,
    coltitle=black,
    boxsep=2pt
]
\begin{tabular}{@{}p{4.4cm}l@{}}
    Architecture         & 2 (input) $\to$ 4 (hidden, sigmoid) $\to$ 1 (output, sigmoid) \\
    Training Algorithm   & Backpropagation with full-batch gradient descent \\
    Learning Rate ($\eta$) & 0.5 \\
    Epochs                & 5000 \\
    Weight Initialization  & $\mathrm{Uniform}(-1,1)$, fixed seed (\texttt{RandomState(0)}) \\
    Final MSE              & 0.000940 \\
    Final Misclassifications & 0 out of 4 \\
\end{tabular}
\end{tcolorbox}

\subsection*{Objective}
Since the single-layer perceptron above was shown to never converge on the XOR
gate, a Multi-Layer Perceptron (MLP) was implemented from scratch in Python to
demonstrate that adding a hidden layer resolves this limitation. As with the
single-layer perceptron, the weights were logged after every update, the
decision boundary was plotted at each logged update, and the results were
analyzed.

\subsection*{Implementation}
A fully-connected MLP with one hidden layer was trained on the XOR truth table
using the forward pass
\begin{equation}
    a_1 = \sigma(X W_1 + b_1), \qquad \hat{y} = \sigma(a_1 W_2 + b_2)
\end{equation}
where $\sigma(z) = 1/(1+e^{-z})$ is the sigmoid activation, followed by
backpropagation to compute gradients and a full-batch gradient descent update
\begin{equation}
    W \leftarrow W + \eta \, \nabla W, \qquad b \leftarrow b + \eta \, \nabla b
\end{equation}
applied once per epoch (i.e.\ one weight update per epoch, unlike the
single-layer perceptron which updated only on misclassified samples). The
network used 2 input neurons, 4 hidden neurons, and 1 output neuron, with
weights initialized from $\mathrm{Uniform}(-1,1)$ using a fixed random seed so
the run is reproducible, and was trained for 5000 epochs with $\eta = 0.5$.
Because a full-batch update happens every epoch, logging all 5000 individual
updates would not be readable, so the weight state was snapshotted every 200
epochs (plus epoch 0, before training).

\subsection*{Weights After Each Update}

\begin{table}[H]
\centering
\small
\caption{MLP Weight Updates on the XOR Gate (snapshotted every 200 epochs; $W_2$ and $b_2$ are the output-layer weights and bias)}
\begin{tabular}{@{}cccccccc@{}}
\toprule
\textbf{Update} & \textbf{Epoch} & \textbf{MSE} & \textbf{Misclassified} & \multicolumn{4}{c}{\textbf{$W_2$}} \\
\midrule
0  & 0    & --       & --  & 0.9273 & $-0.2331$ & 0.5835  & 0.0578 \\
1  & 200  & 0.250085 & 2/4 & 0.7510 & $-0.3583$ & 0.4021  & $-0.0274$ \\
2  & 400  & 0.250020 & 2/4 & 0.7334 & $-0.3496$ & 0.3834  & 0.0274 \\
3  & 600  & 0.249989 & 2/4 & 0.7234 & $-0.3516$ & 0.3723  & 0.0724 \\
4  & 800  & 0.249952 & 3/4 & 0.7150 & $-0.3596$ & 0.3631  & 0.1237 \\
5  & 1000 & 0.249842 & 1/4 & 0.7038 & $-0.3763$ & 0.3507  & 0.2041 \\
6  & 1200 & 0.249233 & 1/4 & 0.6791 & $-0.4192$ & 0.3224  & 0.3784 \\
7  & 1400 & 0.242628 & 1/4 & 0.5945 & $-0.5908$ & 0.2170  & 0.9398 \\
8  & 1600 & 0.200315 & 1/4 & 0.5471 & $-1.2517$ & 0.0020  & 2.4946 \\
9  & 1800 & 0.125264 & \textbf{0/4} & 1.2061 & $-2.9970$ & $-0.1627$ & 4.1044 \\
10 & 2000 & 0.035170 & 0/4 & 1.5213 & $-5.2277$ & $-0.5867$ & 5.6137 \\
11 & 2200 & 0.013730 & 0/4 & 1.4981 & $-6.3138$ & $-0.8090$ & 6.5105 \\
12 & 2400 & 0.007779 & 0/4 & 1.4620 & $-6.9020$ & $-0.9278$ & 7.0316 \\
13 & 2600 & 0.005260 & 0/4 & 1.4335 & $-7.2904$ & $-1.0058$ & 7.3857 \\
14 & 2800 & 0.003916 & 0/4 & 1.4113 & $-7.5765$ & $-1.0631$ & 7.6509 \\
15 & 3000 & 0.003095 & 0/4 & 1.3936 & $-7.8015$ & $-1.1082$ & 7.8617 \\
16 & 3200 & 0.002546 & 0/4 & 1.3791 & $-7.9863$ & $-1.1452$ & 8.0360 \\
17 & 3400 & 0.002156 & 0/4 & 1.3671 & $-8.1427$ & $-1.1765$ & 8.1844 \\
18 & 3600 & 0.001865 & 0/4 & 1.3568 & $-8.2781$ & $-1.2037$ & 8.3134 \\
19 & 3800 & 0.001641 & 0/4 & 1.3480 & $-8.3974$ & $-1.2276$ & 8.4275 \\
20 & 4000 & 0.001463 & 0/4 & 1.3403 & $-8.5038$ & $-1.2490$ & 8.5295 \\
21 & 4200 & 0.001318 & 0/4 & 1.3335 & $-8.5999$ & $-1.2684$ & 8.6219 \\
22 & 4400 & 0.001199 & 0/4 & 1.3275 & $-8.6873$ & $-1.2860$ & 8.7062 \\
23 & 4600 & 0.001099 & 0/4 & 1.3221 & $-8.7676$ & $-1.3022$ & 8.7838 \\
24 & 4800 & 0.001013 & 0/4 & 1.3172 & $-8.8418$ & $-1.3172$ & 8.8555 \\
25 & 5000 & 0.000940 & 0/4 & 1.3128 & $-8.9107$ & $-1.3312$ & 8.9222 \\
\bottomrule
\end{tabular}
\end{table}

\noindent Unlike the single-layer perceptron's table,which went into a cycle and never reached zero misclassifications, the MLP's
misclassification count drops to 0 out of 4 starting at Update 9 (epoch 1800)
and never rises again for the remaining 4000+ epochs, while the MSE continues
to shrink smoothly towards zero.

\subsection*{Decision Boundary Evolution}

\begin{figure}[H]
    \centering
    \includegraphics[width=1\textwidth]{mlp_xor_decision_boundary_evolution.pdf}
    \caption{MLP Decision Boundary After Each of the 26 Logged Weight Updates (Update 0 = initial state, before training)}
    \label{fig:mlp_xor_boundary}
\end{figure}
\noindent\textbf{Inference:} The decision boundary starts out as a roughly
straight, mostly-arbitrary split (epoch 0--600), then bends into a
curved, non-linear region between epoch 800 and 1800, at which it settles into a stable diagonal band that correctly separates the two
$(0,1)/(1,0)$ points from the two $(0,0)/(1,1)$ points. From epoch 2000 onward
the boundary shape stays mostly the same and only sharpens slightly,
confirming convergence.

\begin{figure}[H]
    \centering
    \includegraphics[width=0.8\textwidth]{mlp_xor_training_error.pdf}
    \caption{MLP Trained on XOR -- Convergence}
    \label{fig:mlp_xor_error}
\end{figure}
\noindent\textbf{Inference:} The number of misclassified samples drops from 2
out of 4 at the start, briefly rises to 3 out of 4 around epoch 800 while the
boundary is still adjusting, then falls to 0 out of 4 by epoch 1800 and
stays at 0 for the rest of training. This is the direct opposite of the permanent
error cycle that was seen with the single-layer perceptron.

\begin{figure}[H]
    \centering
    \includegraphics[width=0.85\textwidth]{mlp_xor_training_loss.pdf}
    \caption{MLP Trained on XOR -- Loss Curve (MSE vs Epoch)}
    \label{fig:mlp_xor_loss}
\end{figure}
\noindent\textbf{Inference:} The MSE loss stays close to 0.25 for roughly the first 1000 epochs, then drops sharply between
epoch 1000 and epoch 2200, and finally decays smoothly and asymptotically
towards zero (reaching 0.00094 by epoch 5000). This smooth downward trend,
with no oscillation, confirms a converging model.

\begin{figure}[H]
    \centering
    \includegraphics[width=0.85\textwidth]{mlp_xor_weight_evolution.pdf}
    \caption{MLP Weight Evolution Over Logged Updates}
    \label{fig:mlp_xor_weights}
\end{figure}
\noindent\textbf{Inference:} All four parameter groups ($\|W_1\|$, $\|b_1\|$,
$\|W_2\|$, $\|b_2\|$) stay nearly flat for the first 1000 epochs, then
grow rapidly during the same epoch 1000--2200 window where the loss drops
sharply, and finally flatten out again as training approaches epoch 5000. This shows that the network settles into a stable set of weights rather than oscillating indefinitely,
unlike the single-layer perceptron in Figure~\ref{fig:xor_weights}.

\subsection*{Analysis: Why the MLP Converges on XOR (Where the Single Layer Perceptron Could Not)}
The single-layer perceptron could never converge on XOR because it is
restricted to a single straight-line decision boundary, and \textbf{XOR is not
linearly separable}. The MLP resolves this by introducing a hidden layer: each
of the 4 hidden neurons learns its own straight-line boundary (its own
weighted sum passed through a sigmoid), and the output neuron then combines
these 4 lines into a single non-linear decision region, visible as the curved
boundary in Figure~\ref{fig:mlp_xor_boundary}. The MLP helps in composing several linear boundaries through a hidden layer, and thus produces exactly the kind of non-linear region needed to place all four XOR points on
their correct sides simultaneously. This is the core reason a Multi-Layer
Perceptron is a universal function approximator for problems like XOR that a
single-layer perceptron can never solve.
% ============================================================
% 10. DISCUSSION
% ============================================================
\section{Discussion}

\begin{enumerate}[label=\arabic*.]

\item \textbf{Which optimization method was used?}\\
RandomizedSearchV with 5-fold cross-validation was used, sampling 15
candidate hyperparameter combinations from the full search space rather than
searching and evaluating all of them.

\item \textbf{Which hyperparameters were selected?}\\
The best configuration found was: 1 hidden layer of 128 neurons, sigmoid
activation, no dropout, the Adam optimizer with a learning rate of 0.001,
a batch size of 128, and 30 training epochs.

\item \textbf{Did optimization improve performance?}\\
Yes, though slightly: test accuracy improved from 0.8836 (baseline) to 0.8871
(optimized), with corresponding small gains in precision, recall, and
F1-score. However, the optimized model used a simpler architecture(a single hidden layer) than the two-hidden-layer baseline, yet still performs
slightly better, showing that for this dataset, tuning the learning rate,
batch size, and training length mattered more than adding architectural depth.

\item \textbf{Which hyperparameter had the greatest impact?}\\
Based on the spread seen in the hyperparameter search results plot
(Figure~\ref{fig:search_results}), the learning rate and optimizer choice had
the most visible impact on cross-validation accuracy -- configurations using a
learning rate of 0.1 tended to perform considerably worse, likely due to
overshooting weight updates, while learning rates of 0.001--0.01
paired with Adam or RMSProp consistently gave stronger results.

\item \textbf{Compare Grid Search and Randomized Search.}\\
Grid Search exhaustively evaluates every combination in the search space,
guaranteeing the best combination within the grid is found, but its
cost grows eponentially with the number of hyperparameters and their
candidate values which is infeasible here, since the full grid
($3\times4\times3\times4\times3\times3\times3\times3 = 34{,}992$ combinations)
would require tens of thousands of model fits. Randomized Search instead
samples a fixed number of combinations (15, here) from the same space,which completes in a
practical amount of time, which is why it was the recommended and chosen
method for this experiment.

\item \textbf{Recommend the best MLP configuration.}\\
For this dataset and task, the recommended configuration is the one found by
the search: a single hidden layer of 128 neurons with sigmoid activation, no
dropout, trained with Adam (learning rate 0.001), batch size 128, for 30
epochs. This configuration achieved the highest test accuracy (0.8871) while
also being computationally lighter than the two-hidden-layer baseline.

\end{enumerate}

% ============================================================
% 11. REPORT CONTENTS (self-reference, per manual requirement)
% ============================================================
\section{Report Contents}
This report includes the Objective, Background Theory, Dataset description,
Experimental Procedure (Tasks 1--5), Hyperparameter Optimization methodology
and search space, all mandatory Plots with inference, Results (best
hyperparameters and performance comparison), the Additional Task (K4:
Perceptron Learning Algorithm for the XOR Gate), Discussion, and References.

% ============================================================
% 12. REFERENCES
% ============================================================
\section{References}
\begin{enumerate}[label=\arabic*.]
    \item I. Goodfellow, Y. Bengio and A. Courville, \textit{Deep Learning}, MIT Press, 2016.
    \item C. M. Bishop, \textit{Pattern Recognition and Machine Learning}, Springer, 2006.
    \item S. Haykin, \textit{Neural Networks and Learning Machines}, Pearson, 2009.
    \item Fashion-MNIST Dataset, Zalando Research.
    \item TensorFlow/Keras Documentation.
    \item SciKeras Documentation.
\end{enumerate}

\end{document}